\subsection{Projective schemes}

\blemm

    Suppose we have a collection of schemes $\set{X_i}_{i\in I}$, open subschemes $X_{ij}\subseteq X_i$ for $i,j\in I$
    with $X_{ii}=X_i$, and isomorphisms $f_{ij}\colon X_{ij}\cong X_{ji}$.
    Further suppose that these isomorphisms satisfy the {\emphcolor cocycle condition}:
    $$ f_{ik}|_{X_{ij}\cap X_{ik}} = f_{jk}|_{X_{ji}\cap X_{jk}}\circ f_{ij}|_{X_{ij}\cap X_{ik}} . $$
    In this condition we leave implicit that $f_{ij}(X_{ij}\cap X_{ik})\subseteq X_{jk}$.

    Then there exists a unique scheme $X$ which has open subschemes isomorphic to $X_i$ whose intersections are isomorphic
    to $X_{ij}$, which respects the gluing data in the obvious sense.

\elemm

\bproof

    \TODO{This}

\eproof

\bexam[title={Projective space}]

    Let $A$ be a ring, we define {\emphcolor projective $n$-space} $\bP^n_A$ by gluing together $n+1$ copies of
    $\bA^n_A=\spec A[x_1,\dots,x_n]$ in a particular manner.
    Let us denote these copies by $U_0,\dots,U_n$ where we write $U_i$ as
    $$ U_i = \spec A[x_{0/i},\dots,x_{n/i}]/(x_{i/i}-1) . $$
    Then we let $U_{ij}=\dof(x_{j/i})$, and we identify $U_{ij}$ and $U_{ji}$ via the isomorphism
    $$ A[x_{0/i},\dots,x_{n/i},1/x_{j/i}]/(x_{i/i}-1) \cong A[x_{0/j},\dots,x_{n/j},1/x_{i/j}]/(x_{j/j}-1) $$
    which maps $x_{k/i}$ to $x_{k/j}/x_{i/j}$.
    Because $\spec$ is an equivalence, this forms an isomorphism of schemes.

    It is easy to see that these isomorphisms satisfy the cocycle conditions, so we can glue them together to get
    the scheme $\bP^n_A$.

    We can write $\bP^n_A$ as a colimit of the $U_i$ and $U_{ij}$s, and because $\Gamma$ contravariantly preserves colimits
    (\refmath[corollary]{specsectlim}), we can write $\Gamma(\bP^n_A)$ as a limit of $A[x_{0/i},\dots,x_{n/i}]/(x_{i/i}-1)$
    and $A[x_{0/i},\dots,x_{n/i},1/x_{j/i}]/(x_{i/i}-1)$.
    It is easy to see that this colimit is isomorphic to $A$, so that $\Gamma(\bP^n_A)\cong A$.
    But for $n>0$, $\bP^n_A$ is not isomorphic to $\spec A$, so $\bP^n_A$ is not affine.

\eexam

We can generalize the previous example.

\bdefn

    If $M$ is a commutative monoid, a {\emphcolor $M$-graded ring} is a ring $S_\bul=\bigoplus_{m\in M}S_m$ (where $S_m$ are
    Abelian groups, whose elements are called {\emphcolor homogeneous of degree $m$}) such that multiplication in $S_\bul$
    respects the grading: it forms maps $S_m\times S_n\to S_{n+m}$ for $n,m\in M$.

    Clearly then $S_0$ is a subring and $S_\bul$ is an $S_0$-algebra.
    An ideal $I$ of $S_\bul$ is {\emphcolor homogeneous} if it is generated by homogeneous elements.

\edefn

If $S_\bul$ is $M$-graded, then there are obvious projections $S_\bul\to S_m$ for all $m\in M$; the image of $a\in S_\bul$
under this projection is called its {\emphcolor degree $m$ part} and is denoted $a_m$.
An ideal $I\leq S_\bul$ is homogeneous iff it is closed under these projections: if $a\in I$ then $a_m\in I$.
Clearly if this holds then $I$ is homogeneous: it is generated by the images of itself under all the projections.
Conversely if $I$ is homogeneous, say generated by $I_\homo$, then it is easily seen that any combination of elements in
$I_\homo$ have degree $m$ parts which are also combinations of $I_\homo$.

Thus a homogeneous ideal can be decomposed into homogeneous pieces $I=\bigoplus_{m\in M}I_m$ where $I_m\leq S_m$ is a subgroup
(and $I_0\leq S_0$ is an ideal).
Thus $S_\bul/I\cong\bigoplus_{m\in M}S_m/I_m$ has a natural $M$-grading.

Clearly the sum and product of homogeneous ideals are homogeneous.
The intersection of homogeneous ideals is homogeneous as well: if $a\in I\cap J$ then $a_m\in I\cap J$ for all $m\in M$
so by the above consideration $I\cap J$ is homogeneous.
And if $I$ is homogeneous, so is $\sqrt I$: if $a^n\in I$ then
$$ (a^n)_m = \sum_{k_1+\cdots+k_n=m}a_{k_1}\cdots a_{k_n} $$
is in $I$ as all the $a_{k_i}$ are.

If $M$ is a commutative monoid, we associate to it the Abelian group $\tilde M$ which is the quotient of the monoid
$M\times M$ by the equivalence $(m,n)\sim(m+k,n+k)$ for all $k\in M$.
We can embed $M\to\tilde M$ by $m\mapsto(m,0)$, and we note that this is the identity of an adjunction
$\CommMon\adjarrow{}{}\Ab$.

If $T\subseteq S_\bul$ is a multiplicative subset containing only homogeneous elements, then $T^{-1}S_\bul$ has an obvious
grading over $\tilde M$ where $a_m/t_n$ has degree $m-n$.
It is easy to check that this a well-defined grading.

\bnote

    In order to avoid excessive subscripts, we will also use the notation $S_\bul[T^{-1}]$ for $T^{-1}S_\bul$.
    In particular, $(S_\bul)_f$ will instead be denoted $S_\bul[1/f]$.

\enote

Now we consider only the case where $M=\bZ$ or $\bZ_{\geq0}$ (every $\bZ_{\geq0}$-grading can be turned into a $\bZ$ grading
by setting $S_n=0$ for $n<0$).
We define the {\emphcolor degree} of $a\in S_\bul$ to be
$$ \deg(a) = \max\set{n\in\bZ}[a_n\neq0] . $$

In order for an homogeneous ideal $I<S_\bul$ to be prime, it is sufficient to check the usual condition on homogeneous elements
only: that is, if $ab\in I$ then $a,b\in I$ for homogeneous $a,b$.
Suppose that $a,b\notin I$, then define $a_I=\sum_{m,a_m\in I}a_m$ and $a^I=\sum_{m,a_m\notin I}a_m$ so that $a=a_I+a^I$, thus
$$ ab = a_Ib_I + a_Ib^I + a^Ib_I + a^Ib^I . $$
As $a_I\in I$ we have that $a^I\notin I$, and we need to show that $a^Ib^I\notin I$.
That is, it is sufficient to consider the case where $a_m\notin I$ for all $m\in M$ (or is zero).
Let $n=\deg(a)$ and $m=\deg(b)$ so that $(ab)_{n+m}=a_nb_m$.
Because $a_n,b_m\notin I$ we have that $(ab)_{n+m}\notin I$ (by primality on homogeneous elements) and thus $ab\notin I$
(because $I$ is homogeneous).

\bnote

    For the remainder of these notes, $S_\bul$ will be a $\bZ_{\geq0}$-graded ring unless otherwise specified.

\enote

If $S_0=A$, then $S_\bul$ is said to be {\emphcolor graded over $A$}.
For example $A[x_0,\dots,x_n]$ is graded over $A$ (where the homogeneous elements are homogeneous polynomials), or $A[x_0,\dots,x_n]/I$ where $I$ is
an homogeneous ideal and $I\cap A=0$.
For a graded ring $S_\bul$, $S_+=\bigoplus_{n>0}S_n$ is an ideal called the {\emphcolor irrelevant ideal}.
If $S_+$ is a finitely generated ideal, then $S_\bul$ is said to be {\emphcolor finitely generated over $A$}.

Note that $S_\bul$ is finitely generated over $A$ iff $S_\bul$ is a finitely generated graded $A$-algebra (i.e.\ generated by finitely many homogeneous elements).
If $S_\bul$ is finitely generated over $A$, so $S_+=(f_1,\dots,f_n)$ homogeneous, then by considering degrees we see that $S_1$ is $A$-spanned by the $f_i$.
Then we can see that $S_2$ is $A$-spanned by the $f_i,f_if_j$s, and so on by induction.
That is, $S_k$ is spanned by the products of $f_i$s, so that the $f_i$s generate $S_\bul$ as an $A$-algebra.
Conversely if $S_\bul$ is finitely generated as a graded $A$-algebra by $f_1,\dots,f_n$, then these must also generate $S_+$.

Furthermore $S_\bul$ is Noetherian iff $A$ is, and $S_\bul$ is finitely generated over $A$.
Indeed if $S_\bul$ is Noetherian, then $S_+$ is finitely generated as it is an ideal; by the previous paragraph this means that $S_\bul$ is finitely generated.
And any ideal of $A$ can be extended to an homogeneous ideal of $S_\bul$ (with the same generators), which is finitely generated.
It is clear that the degree zero generators of this ideal generate the ideal from $A$.
Alternatively we could note that the extension of $A$'s ideal is unique, as it adds no additional elements from $A$, and consider the definition of a Noetherian ring.
Conversely if $S_\bul$ is finitely generated and $A$ is Noetherian, then $S_\bul$ is a Noetherian $A$-algebra and thus a Noetherian ring.

Now we can see why graded rings are important.
The idea is that for a graded ring $S_\bul$ we will define a scheme $\Proj S_\bul$, motivated by $\bP^n_A$.
That is, $\Proj A[x_0,\dots,x_n]$ will give us $\bP^n_A$.

\bprop

    Let $f\in S_+$ be homogeneous.
    Then there is a bijection between the prime ideals of $(S_\bul)[1/f]_0$ and the homogeneous prime ideals of $(S_\bul)[1/f]$.

\eprop

\bproof

    We generalize this result as follows: for any $\bZ$-graded ring $A$ with an invertible homogeneous element $f\in A$ of positive degree,
    there is a bijection between the prime ideals of $A_0$ and the homogeneous prime ideals of $A$.

    Consider the ring map $A_0\to A$, which then maps homogeneous prime ideals of $A$ to prime ideals of $A_0$ via the inverse image.
    Note that this maps a prime ideal $P=\bigoplus_iQ_i\leq A$ to $Q_0$.
    Conversely if $P_0\leq A_0$ is a prime ideal, then define $P\leq A$ as $P=\bigoplus_iQ_i$ where $a\in Q_i$ iff $a^{\deg f}/f^i\in P_0$.
    Here we have that $Q_0=P_0$.
    It remains to show that $P$ is an ideal, and prime.

    Note that $a\in Q_i$ iff $a^2\in Q_{2i}$: $a^{\deg f}/f^i\in P_0$ iff $a^{2\deg f}/f^{2i}\in P_0$ as $P_0$ is prime, iff $a^2\in Q_{2i}$.
    And if $a_1,a_2\in Q_i$ then $(a_1+a_2)^2\in Q_{2i}$; indeed by the binomial theorem (letting $d=\deg f$)
    $$ \frac{(a_1+a_2)^{2d}}{f^{2i}} = \frac{\sum_k\binom{2d}ka_1^ka_2^{2d-k}}{f^{2i}} , $$
    and for all $k$, one of $k$ or $2d-k$ is at least $d$, so each summand is in $P_0$.
    Thus $a_1+a_2\in Q_i$, meaning $P$ is closed under addition.

    And if $a\in Q_i$ and $b\in S_\bul$ has degree $k$, then $ab\in Q_{i+k}$ because
    $$ \frac{(ab)^d}{f^{i+k}} = \frac{a^d}{f^i}\cdot\frac{b^d}{f^k} \in P_0 . $$
    Thus $P$ is closed under multiplication so it is an ideal.

    It is sufficient to check primality on homogeneous elements, so if $a,b\in S_\bul$ are homogeneous of degrees $i,j$ respectively,
    then if $ab\in Q_{i+j}$ that means $(ab)^d/f^{i+j}\in P_0$.
    By primality $a^d/f^i,b^d/f^j\in P_0$ so $a\in Q_i$ and $b\in Q_j$ as desired.
    Thus $P\leq S_\bul$ is a prime ideal.

    Now if $P\leq S_\bul$ is a prime ideal with $P=\bigoplus_iP_i$, we want to show that $a\in P_i$ iff $a^d/f^i\in P_0$.
    Indeed if $a^d/f^i\in P_0$ then by primality $a\in P_i$.
    Conversely if $a\in P_i$ then $a^d/f^i\in P$ is homogeneous of degree $0$, so it is in $P_0$ as desired.
    So the two maps are mutual inverses, as desired.
    \qed

\eproof

Now we can construct $\Proj S_\bul$ as the gluing of $\spec S_\bul[1/f]_0$ along the $U_{fg}=\dof(g^{\deg f}/f^{\deg g})$, where $f$ ranges over the homogeneous elements
of positive degree.
Note that for $S_\bul=A[x_0,\dots,x_n]$ and $f=x_i$ then $S_\bul[1/f]_0\cong A[x_{0/i},\dots,x_{n/i}]/(x_{i/i}-1)$ by $x_{k/i}\mapsto x_k/x_i$.
And $U_{ij}=\dof(x_j/x_i)$ corresponds to $\dof(x_{j/i})$, so this gives us the original definition of $\bP^n_A$ (we will see that we need only glue along generators).

We note that $U_{fg}$ is isomorphic to $\spec S_\bul[1/fg]_0$.
This isomorphism is induced from the map $S_\bul[1/f]\to S_\bul[1/fg]$.
Thus there is an isomorphism $U_{fg}\cong U_{gf}$, and it is routine to check that these isomorphisms satisfy the cocycle condition.
Thus we have defined $\Proj S_\bul$; we will continue this section by providing more descriptions of $\Proj S_\bul$.

By the above proposition, the prime ideals of $S_\bul[1/f]_0$ correspond to the homogeneous prime ideals of $S_\bul[1/f]$.
We can construct a bijection between the points of $\Proj S_\bul$ and the homogeneous prime ideals of $S_\bul$.
For a homogeneous prime ideal $\p\leq S_\bul$, then for all $f\notin\p$ we have the homogeneous prime ideal $\p_f\in S_\bul[1/f]$.
All of these ideals are identified: if $f,g\notin\p$ then $fg\notin\p$ so that $\p_f$ and $\p_g$ are identified in $\spec S_\bul[1/fg]_0$.
And all of homogeneous prime ideals of $S_\bul[1/f]$ are of this form, so mapping $\p\mapsto\p_f$ is surjective.
Because $\p\mapsto\p_f$ is an injection (for fixed $f$, $f\notin\p$), we see that distinct ideals are not identified.

Note that we must consider only $\p$ which don't contain some homogeneous $f\in S_+$; that is $S_+$ cannot be a subset of $\p$.
This is because we consider only $f$ with positive degree.
These homogeneous prime ideals which do not contain $S_+$ are called {\emphcolor relevant}; otherwise they are {\emphcolor irrelevant}.

Thus the points of $\Proj S_\bul$ correspond to the relevant homogeneous prime ideals of $S_\bul$.

\bdefn

    Let $T\subseteq S_\bul$ be a collection of homogeneous elements.
    Define its {\emphcolor vanishing set} to be
    $$ \vof(T) = \set{\p\leq S_\bul\hbox{ homogeneous prime}}[T\subseteq\p,\hbox{ and $\p$ does not contain the irrelevant ideal}] . $$
    And if $f$ is a homogeneous element of {\it positive degree\/}, define $\dof_+(f)=\Proj S_\bul-\vof(f)$; its {\emphcolor projective distinguished open set}.

\edefn

Note that $\dof_+(f)$ can be written as
$$ \dof_+(f) = \set{\p\leq S_\bul\hbox{ homogeneous prime}}[f\notin\p] . $$
It is immediat that this does not contain any irrelevant ideals as $f\in S_+$.
Furthermore this corresponds to $\spec S_\bul[1/f]_0$.
Indeed, the homogeneous prime ideals of $S_\bul[1/f]$ correspond to the homogeneous prime ideals of $S_\bul$ not containing $f$.

The projective vanishing sets satisfy the same properties as those in the affine case.
Thus they form a topology of closed sets on $\Proj S_\bul$, and $\dof_+(f)$ form a basis.

Note that any homogeneous ideal $I\leq S_\bul$ is contained in a homogeneous prime ideal.
Indeed, for an ideal $I$ define its {\emphcolor homogenization} to be
$$ I_\homo = \bigoplus_n(I\cap S_n) = \ideal{a\in I}[\hbox{$a$ is homogeneous}] . $$
Because primality can be checked on homogeneous elements, if $\p$ is prime so is $\p_\homo$.
Thus if $I\leq\p$ is prime, then $I\leq\p_\homo$ is homogeneous prime.
Note that because $\p_\homo\leq\p$, we have that
$$ \sqrt I = \bigcap_{I\leq\p\hbox{ homogeneous}}\p . $$
Here the intersection is almost over $\vof(I)$, but it may be the case that the $\p$s are irrelevant.
We will remedy this now.

Note that the following are equivalent for any homogeneous ideal $I\subseteq S_+$.
\benum
    \item $\vof(I)=\emptyset$,
    \item if $I=(f_i)_i$ are homogeneous generators, then $\bigcup_i\dof_+(f_i)=\Proj S_\bul$,
    \item $\sqrt I\supseteq S_+$.
\eenum
(1) and (2) are equivalent because
$$ \vof(f_i)_i = \bigcap_i\vof(f_i) , $$
which is then empty iff its complement is $\Proj S_\bul$, which is precisely (2).
Note that $\vof(I)=\emptyset$ iff all of the homogeneous prime ideals containing $I$ also contain $S_+$.
This is iff their intersection contains $S_+$.
Because this intersection is $\sqrt I$, this gives us (3).

Call an ideal $I$ satisfying any of the above equivalent conditions {\emphcolor irrelevant}.
Clearly by (3) this agrees with our definition of an irrelevant prime ideal.

For a subset $Z\subseteq\Proj S_\bul$, define the relevant homogeneous ideal $\idof(Z)\leq S_\bul$ by
$$ \idof(Z) = \bigcap_{\p\in Z}\p . $$
The intersection of homogeneous ideals is homogeneous, so $\idof(Z)$ is homogeneous.
Then we have that
$$ \vof\idof(Z) = \cl Z $$
whose proof is the same as in the affine case.

We also have that for $I\leq S_+$ relevant homogeneous,
$$ \idof\vof(I) = \sqrt I . $$
Indeed because $I\leq S_+$ we can ignore degree zero, in which case all irrelevant prime ideals all look like $S_+$, and thus contribute
nothing to the intersection, as there exists a relevant prime ideal.
Thus the intersection of all homogeneous prime ideals containing $I$ looks the same as the intersection of all relevant prime ideals
containing $I$.

Thus as in the affine case, this clearly demonstrates that $\vof(I)\subseteq\vof(J)$ iff $\sqrt J\leq\sqrt I$.

Now we want to define the structure sheaf on $\Proj S_\bul$.
Recall that there is a correspondence $\phi\colon\dof_+(f)\cong\spec S_\bul[1/f]_0$; we claim that this correspondence is a homeomorphism.
Indeed, the map $\p\mapsto(\p_f)_0$ is continuous: for $g/f^k\in S_\bul[1/f]_0$ we have that
$$ \phi^{-1}\dof(g/f^k) = \set{\p\in\dof_+(f)}[g/f^k\notin(\p_f)_0] . $$
This is equal to $\set{\p}[g\notin\p]$; because $f\notin\p$, we have that $g/f^k\in\p_f$ iff $gf^t\in\p$ for some $t$ iff $g\in\p$.
Thus
$$ \phi^{-1}\dof(g/f^k) = \dof_+(g) $$
so that $\phi$ is continuous.
And it is open, since for $g$ homogeneous
$$ \phi\dof_+(fg) = \set{(\p_f)_0}[fg\notin\p] = \dof\parens{\frac{g^{\deg f}}{f^{\deg g}}} . $$
The last equality is because $g^d/f^k\notin\p_f$ iff $g^d\notin\p$ iff $g\notin\p$ when $f\notin\p$.

Now $\dof_+(fg)$ corresponds to $\dof(g^{\deg f}/f^{\deg g})$ in $\spec S_\bul[1/f]_0$ and to $\dof(f^{\deg g}/g^{\deg f})$ in
$\spec S_\bul[1/g]_0$.
We want to define the structure sheaf on $D_+(f)\subseteq\Proj S_\bul$ to be the structure sheaf of $\spec S_\bul[1/f]_0$.
It is routine to check that this agrees with triple intersections (the cocycle condition).

Note that if $S_+=(f_i)_i$ then we can consider only the $\dof_+(f_i)$.
Indeed, this comes down to the fact that if $X=\bigcup_iU_i$ then any sheaf $(X,\O_X)$ can be written as the gluing of the $(U_i,\O_{U_i})$.

So in terms of $A[x_0,\dots,x_n]$, we have that the $x_i$ generate the irrelevant ideal, so we can focus on the structure sheaf restricted to
$$ \dof_+(x_i) \cong \spec A[x_0,\dots,x_n,1/x_i]_0 \cong \spec A[x_{0/i},\dots,x_{n/i}]/(x_{i/i}-1) . $$
Thus defining $\bP^n_A=\spec A[x_0,\dots,x_n]$ agrees with our previous definition.

